\documentclass{article}
\usepackage{amsmath}
\usepackage{amsfonts}

\begin{document}

\section*{Uitwerking Multidimensional Gaussian Normalisation}

Gegeven is de kansdichtheidsfunctie van een \textit{multidimensional Gaussian distribution}:
\begin{equation}
    f_X(\mathbf{x}) = Z \exp\left[ -\frac{1}{2} (\mathbf{x} - \boldsymbol{\mu})^T A (\mathbf{x} - \boldsymbol{\mu}) \right]
\end{equation}
waarbij $A \in \mathbb{R}^{n \times n}$ een symmetrische en \textit{positive definite} matrix is. De normalisatie-eis luidt:
\begin{equation}
    \int_{\mathbb{R}^n} f_X(\mathbf{x}) \, \mathrm{d}^n\mathbf{x} = 1
\end{equation}

\subsection*{Co\"ordinatentransformatie via LDU-decompositie}
Omdat $A$ symmetrisch is ($A = A^T$), schrijven we de LDU-decompositie als $A = L D L^T$, waarbij $L$ een \textit{lower triangular matrix} is met enen op de diagonaal, en $D$ een \textit{diagonal matrix} met elementen $D^i_{\ i} = d_i > 0$.

Definieer de nieuwe integratievariabele:
\begin{equation}
    \mathbf{y} = L^T (\mathbf{x} - \boldsymbol{\mu}) \implies \mathbf{x} = \boldsymbol{\mu} + (L^T)^{-1}\mathbf{y}
\end{equation}

De \textit{Jacobian matrix} $J$ van deze transformatie ($\frac{\partial \mathbf{x}}{\partial \mathbf{y}}$) is:
\begin{equation}
    J = (L^T)^{-1}
\end{equation}
Omdat $L$ \textit{lower-diagonal} is met een eenheidsdiagonaal, is $\det(L) = 1$. Bijgevolg is $\det(L^T) = 1$ en $\det(J) = \det((L^T)^{-1}) = 1$.

\subsection*{Substitutie in de exponent}
De kwadratische vorm in de exponent herschrijft zich algebra\"isch als volgt:
\begin{align}
    (\mathbf{x} - \boldsymbol{\mu})^T A (\mathbf{x} - \boldsymbol{\mu}) &= \left( (L^T)^{-1}\mathbf{y} \right)^T (L D L^T) \left( (L^T)^{-1}\mathbf{y} \right) \\
    &= \mathbf{y}^T (L^T)^{-T} L D L^T (L^T)^{-1} \mathbf{y}
\end{align}
Aangezien de transpositie van een inverse gelijk is aan de inverse van een transpositie, geldt $(L^T)^{-T} = (L^{-1})^{-1} = L^{-1}$. Dit reduceert de vergelijking tot:
\begin{equation}
    \mathbf{y}^T (L^{-1} L) D (L^T (L^T)^{-1}) \mathbf{y} = \mathbf{y}^T D \mathbf{y}
\end{equation}

\subsection*{Berekening van de normalisatiefactor Z}
De $n$-dimensionale integraal valt nu door de diagonaalmatrix $D$ uiteen in $n$ onafhankelijke 1D Gaussiaanse integralen:
\begin{equation}
    \int_{\mathbb{R}^n} f_X(\mathbf{x}) \, \mathrm{d}^n\mathbf{x} = Z \prod_{i=1}^n \left( \int_{-\infty}^{+\infty} \exp\left(-\frac{1}{2} d_i (y^i)^2\right) \, \mathrm{d}y^i \right) = 1
\end{equation}
Met behulp van de standaardoplossing $\int_{-\infty}^{\infty} e^{-a z^2} \mathrm{d}z = \sqrt{\frac{\pi}{a}}$ (waarbij $a = d_i / 2$), bekomen we:
\begin{equation}
    1 = Z \prod_{i=1}^n \sqrt{\frac{2\pi}{d_i}} \implies Z = \prod_{i=1}^n \sqrt{\frac{d_i}{2\pi}}
\end{equation}
Omdat $\det(A) = \det(L D L^T) = \det(L)\det(D)\det(L^T) = 1 \cdot \left(\prod_{i=1}^n d_i\right) \cdot 1$, geldt $\prod_{i=1}^n d_i = \det(A)$. We herschrijven $Z$ definitief als:
\begin{equation}
    Z = \sqrt{\frac{\det(D)}{(2\pi)^n}} = \sqrt{\frac{\det(A)}{(2\pi)^n}}
\end{equation}

\end{document}